\bmtchapitre{Calcul vectoriel et produit scalaire}{Décomposer un vecteur dans une base et utiliser les expressions du produit scalaire pour démontrer des relations métriques.}{Géométrie}
\label{t11s-c13}
Dans ce chapitre, les calculs de longueurs, d’angles et d’aires à partir de coordonnées se font dans un repère orthonormé.
\begin{revision}Chasles, coordonnées, Pythagore, angles et trigonométrie; repère orthonormé pour les calculs métriques.\end{revision}
\begin{automatismes}\exo\label{t11s-c13-auto}Pour $A(1;2),B(4;6)$, calculer les coordonnées de $\vect{AB}$ et $AB$.\end{automatismes}
\begin{corrige}[exercice~\ref{t11s-c13-auto}]\label{sol-t11s-c13-auto}$\vect{AB}=(3;4)$, $AB=\sqrt{3^2+4^2}=5$.\end{corrige}

\section{Combinaisons linéaires et bases}
\begin{definition}Une combinaison linéaire de $\vecu,\vecv$ est $a\vecu+b\vecv$ avec $a,b\in\R$. Deux vecteurs non colinéaires forment une base du plan : chaque vecteur s’écrit d’une manière unique $a\vecu+b\vecv$.\end{definition}
\begin{propriete}Dans un repère, $\vecu=(x;y)$ et $\vecv=(x';y')$ sont non colinéaires si et seulement si $xy'-yx'\ne0$.\end{propriete}
\begin{demonstration}Le système $a\vecu+b\vecv=\vecw$ a deux équations. L’élimination fournit une solution unique si le déterminant $xy'-yx'$ est non nul. S’il est nul, les colonnes sont proportionnelles (ou l’un des vecteurs est nul), donc colinéaires.\end{demonstration}
\begin{exemple}$\vecu=(1;1)$, $\vecv=(1;-1)$ forment une base. Pour $\vecw=(3;1)$, $a+b=3$ et $a-b=1$, donc $\vecw=2\vecu+\vecv$. Cette base n’est pas un repère orthonormé : ses deux vecteurs ont norme $\sqrt2$.\end{exemple}
\section{Quatre expressions d’un même produit}
\begin{definition}Pour deux vecteurs non nuls, $\vecu\cdot\vecv=\norme{\vecu}\norme{\vecv}\cos\theta$, où $\theta$ est l’angle entre eux. Si l’un est nul, le produit scalaire est nul.\end{definition}
\begin{propriete}
Dans un repère orthonormé, $\vecu\cdot\vecv=xx'+yy'$. Pour $O\ne A$, avec $H$ projeté orthogonal de $B$ sur la droite $(OA)$ et $\vecu=\vect{OA},\vecv=\vect{OB}$, le produit vaut $OA$ multiplié par la longueur algébrique $OH$ orientée vers $A$. Enfin,
\[
\vecu\cdot\vecv=\frac12\big(\norme{\vecu+\vecv}^2-\norme{\vecu}^2-\norme{\vecv}^2\big).
\]
Le produit est symétrique et bilinéaire; $\vecu\cdot\vecu=\norme{\vecu}^2$; deux vecteurs sont orthogonaux si et seulement si leur produit est nul (y compris le vecteur nul).
\end{propriete}
\begin{demonstration}Décomposer dans deux directions orthonormées donne $xx'+yy'$. Développer $(x+x')^2+(y+y')^2$ donne l’identité par les normes. La projection donne $OH=OB\cos\theta$ avec son signe, donc la même expression géométrique. La bilinéarité résulte de la formule des coordonnées.\end{demonstration}
\section{Relations métriques dans un triangle}
\begin{propriete}[Al-Kashi et médiane]Dans un triangle $ABC$,
\[
BC^2=AB^2+AC^2-2AB\,AC\cos\widehat{BAC}.
\]
Si $I$ est milieu de $[BC]$, alors $AB^2+AC^2=2AI^2+BC^2/2$.
\end{propriete}
\begin{demonstration}Al-Kashi résulte du carré de $\vect{BC}=\vect{AC}-\vect{AB}$. Pour la médiane, $\vect{AB}=\vect{AI}+\vect{IB}$ et $\vect{AC}=\vect{AI}-\vect{IB}$; en additionnant les carrés, les produits croisés s’annulent.\end{demonstration}
\begin{propriete}[aire et sinus]Pour un triangle non aplati, $\mathcal A=\dfrac12 AB\,AC\sin\widehat{BAC}$ et
\[
\frac{BC}{\sin\widehat{BAC}}=
\frac{CA}{\sin\widehat{ABC}}=
\frac{AB}{\sin\widehat{ACB}}.
\]
\end{propriete}
\begin{demonstration}La hauteur relative à $AB$ vaut $AC\sin\widehat{BAC}$. Écrire l’aire avec chacun des trois angles puis diviser par les produits des côtés (tous non nuls) donne la relation des sinus.\end{demonstration}

% ENRICHISSEMENT C13

\subsection*{Choisir la bonne expression du produit scalaire}
\begin{methode}{relier angle, projection et coordonnées}
Avec des coordonnées dans un repère orthonormé, calculer d’abord les vecteurs puis leur produit. Pour un angle, diviser par les deux normes non nulles. Pour projeter $C$ sur la droite $(AB)$, $A\ne B$, poser $H=A+t\vect{AB}$; la condition $\vect{HC}\cdot\vect{AB}=0$ impose
\[
t=\frac{\vect{AC}\cdot\vect{AB}}{AB^2}.
\]
Si $t\notin[0;1]$, le pied appartient à la droite mais pas au segment; la hauteur d’un triangle peut être extérieure.
\end{methode}
\begin{exemple}[une projection calculée et vérifiée]
Pour $A(1;1)$, $B(5;3)$, $C(2;4)$, $\vect{AB}=(4;2)$, $\vect{AC}=(1;3)$. Le produit vaut $10$ et $AB^2=20$, donc $t=1/2$, $H(3;2)$. On vérifie $\vect{HC}=(-1;2)$ et son produit avec $(4;2)$ vaut $0$. La distance $CH=\sqrt5$ est la hauteur relative à $AB=2\sqrt5$, d’où l’aire $5$ unités d’aire. Le même produit donne $\cos\widehat{BAC}=10/(\sqrt{20}\sqrt{10})=1/\sqrt2$ : l’angle vaut $\pi/4$.
\end{exemple}
\begin{methode}{contrôler la formule dans une base quelconque}
Les coordonnées d’un vecteur sont les coefficients de sa décomposition dans la base choisie. Si la base n’est pas orthonormée, les produits des vecteurs de base interviennent. On peut revenir aux coordonnées dans un repère orthonormé pour éviter une formule mal appliquée. Une base sert à représenter les vecteurs; elle ne garantit pas des unités de longueur égales.
\end{methode}

\section{Exercices et problèmes}
\exo[Une base oblique]\label{t11s-c13-e01}
Écrire $(5;1)$ dans la base $\vecu=(2;1)$, $\vecv=(1;-1)$.
\begin{corrige}[exercice~\ref{t11s-c13-e01}]\label{sol-t11s-c13-e01}
$2a+b=5$, $a-b=1$; $a=2,b=1$. Déterminant $-3\ne0$, donc base et unicité.
\end{corrige}
\exo[Produit scalaire]\label{t11s-c13-e02}
Calculer le produit de $(2;-1)$ et $(3;4)$, leurs normes et le cosinus de leur angle.
\begin{corrige}[exercice~\ref{t11s-c13-e02}]\label{sol-t11s-c13-e02}
Produit $6-4=2$; normes $\sqrt5$ et $5$; cosinus $2/(5\sqrt5)$. Les deux vecteurs sont non nuls.
\end{corrige}
\exo[Triangle rectangle]\label{t11s-c13-e03}
Pour $A(0;0),B(4;0),C(1;\sqrt3)$, déterminer l’angle en $A$ puis calculer $BC$.
\begin{corrige}[exercice~\ref{t11s-c13-e03}]\label{sol-t11s-c13-e03}
$AB=4$, $AC=2$, produit $4$. Cosinus $4/(4\times2)=1/2$, donc angle $\pi/3$. $BC^2=16+4-2\times4\times2\times1/2=12$, donc $BC=2\sqrt3$.
\end{corrige}
\exo[Médiane]\label{t11s-c13-e04}
Dans un triangle, $AB=5$, $AC=7$, $BC=8$. Calculer la longueur de la médiane issue de $A$.
\begin{corrige}[exercice~\ref{t11s-c13-e04}]\label{sol-t11s-c13-e04}
$25+49=2AI^2+64/2$, donc $AI^2=21$, $AI=\sqrt{21}$. Les longueurs satisfont les inégalités triangulaires.
\end{corrige}
\exo[Une formule mal utilisée]\label{t11s-c13-e05}
Dans la base $\vecu=(1;1)$, $\vecv=(1;-1)$, peut-on calculer le produit scalaire de deux vecteurs en multipliant directement leurs coordonnées dans cette base ?
\begin{corrige}[exercice~\ref{t11s-c13-e05}]\label{sol-t11s-c13-e05}
Pas avec la formule sans coefficients : $\vecu$ a coordonnées $(1;0)$ dans cette base, ce qui donnerait $1$ pour sa norme au carré au lieu de $2$. Il faut une base orthonormée, ou utiliser les produits des vecteurs de base.
\end{corrige}
\exo[Mesurer une aire]\label{t11s-c13-e06}
Deux côtés d’une parcelle triangulaire mesurent $30$ m et $40$ m, avec angle compris $60\degre$. Calculer sa surface et son troisième côté.
\begin{corrige}[exercice~\ref{t11s-c13-e06}]\label{sol-t11s-c13-e06}
Aire $\tfrac12\times30\times40\times\sqrt3/2=300\sqrt3\ \mathrm{m}^2$. Troisième côté au carré $900+1600-1200=1300$, longueur $10\sqrt{13}$ m.
\end{corrige}

\exo[Une hauteur hors du segment]\label{t11s-c13-p01}
Dans un repère orthonormé, $A(0;0)$, $B(4;0)$, $C(-2;3)$. Déterminer le projeté orthogonal $H$ de $C$ sur la droite $(AB)$ et dire s’il appartient au segment $[AB]$. Calculer l’aire du triangle et le cosinus de l’angle en $A$; cet angle est-il aigu, droit ou obtus ?
\begin{corrige}[exercice~\ref{t11s-c13-p01}]\label{sol-t11s-c13-p01}
La droite $(AB)$ est $y=0$, donc $H=(-2;0)$. Son abscisse n’est pas dans $[0;4]$ : pied extérieur. La hauteur $CH=3$ donne l’aire $4\times3/2=6$. Le produit $\vect{AB}\cdot\vect{AC}=4(-2)=-8$, avec normes $4$ et $\sqrt{13}$, donne $\cos\widehat{BAC}=-2/\sqrt{13}<0$. L’angle du triangle est donc obtus. Une hauteur extérieure n’empêche pas la formule base fois hauteur.
\end{corrige}
\exo[Des coordonnées dans une base oblique]\label{t11s-c13-p02}
Dans un repère orthonormé de référence, $\vecu=(1;0)$ et $\vecv=(1;1)$. On définit $\vecw=\vecu+\vecv$ et $\vecz=\vecu-\vecv$. Un élève utilise leurs coordonnées $(1;1)$ et $(1;-1)$ dans la base $(\vecu,\vecv)$ pour conclure qu’ils sont orthogonaux. Contrôler que cette paire est une base, puis calculer le vrai produit scalaire et expliquer l’erreur.
\begin{corrige}[exercice~\ref{t11s-c13-p02}]\label{sol-t11s-c13-p02}
Le déterminant de $\vecu,\vecv$ est $1$, donc ils forment une base. Dans le repère orthonormé de référence, $\vecw=(2;1)$ et $\vecz=(0;-1)$ : produit $-1\ne0$, donc ils ne sont pas orthogonaux. La base donnée est oblique : $\vecu\cdot\vecv=1\ne0$. Multiplier directement les coefficients dans cette base revient à supposer, à tort, des vecteurs unitaires orthogonaux.
\end{corrige}

\begin{jemeteste}Je distingue une base quelconque d’un repère orthonormé; je choisis coordonnées, normes, angle ou projection selon les données.\end{jemeteste}
\begin{notepourlaclasse}Le produit scalaire de vecteurs de coordonnées dans une base oblique est un piège utile. Faire établir les relations métriques avec Chasles et bilinéarité.\end{notepourlaclasse}
